\chapter{Green Tensor}

\section{Evaluation of the Imaginary part of the Green tensor at emitter position, $\Im\left\{\mathbf{G}\left(\mathbf{r}_{0},\mathbf{r}_{0}\right)\right\}$\label{App:green_at_origin}}

The Green tensor is used to propagate the electric field generated by a point source at $\mathbf{r}_{0}$ to some point $\mathbf{r}$, represented as $\mathbf{G}\left(\mathbf{r},\mathbf{r}_{0}\right)$. To evaluate the radiated power by an emitter is necessary to calculate the imaginary part of the Green tensor at the emitter position, \textit{i.e.}, $\Im\left\{\mathbf{G}\left(\mathbf{r}_{0},\mathbf{r}_{0}\right)\right\}$. Although its evaluation seams impossible because appears to be infinite for $R = 0$, the problem can be circumvented. The key is in the expansion of the exponential term.\\
The Green tensor reads:

\begin{equation}
\mathbf{G}\left(\mathbf{r},\mathbf{r}_{0}\right) = \frac{e^{ikR}}{4\pi R}\left[\left(1+\frac{ikR-1}{k^2R^2}\right)\mathbb{I}+\frac{3-3ikR-k^2R^2}{k^2R^2}\frac{\mathbf{R} \otimes\mathbf{R}}{R^2}\right],
\end{equation}
with $R = \mathbf{r}-\mathbf{r}_{0}$ and $\frac{\mathbf{R} \otimes\mathbf{R}}{R^2} = \hat{\mathbf{R}}\otimes \hat{\mathbf{R}}$.
The imaginary part of the Green tensor for the case where $\mathbf{r} = \mathbf{r}_{0}$ is:

\begin{eqnarray*}
\Im\left\{ \mathbf{G}\left(\mathbf{r}_{0},\mathbf{r}_{0}\right)\right\}  & = & \lim_{R \rightarrow 0 }\Im\left\{ \frac{e^{ikR}}{4\pi R}\left[\left(1+\frac{ikR-1}{k^{2}R^{2}}\right)\mathbb{I}+\frac{3-3ikR-k^{2}R^{2}}{k^{2}R^{2}}\hat{\mathbf{R}}\otimes \hat{\mathbf{R}}\right]\right\} \\
\end{eqnarray*}
\begin{eqnarray*}
 & = & \lim_{R \rightarrow 0 }\Im\left\{ \frac{1}{4\pi R}\left[1+ikR-\frac{1}{2}\left(kR\right)^{2}-i\frac{1}{6}\left(kR\right)^{3}+\mathcal{O}^{4}\right]\right. \\
 & & \left. \left[\left(1+\frac{ikR-1}{k^{2}R^{2}}\right)\mathbb{I}+\frac{3-3ikR-k^{2}R^{2}}{k^{2}R^{2}}\hat{\mathbf{R}}\otimes \hat{\mathbf{R}}\right]\right\} \\
 & = & \lim_{R \rightarrow 0 } \left\{\frac{1}{4\pi R}\left[\left(\frac{ikR}{k^{2}R^{2}}+ikR-\frac{ikR}{k^{2}R^{2}}-\frac{ik^{3}R^{3}}{2k^{2}R^{2}}-i\frac{1}{6}k^{3}R^{3}+i\frac{k^{3}R^{3}}{6k^{2}R^{2}}\right)\mathbb{I}+ \right. \right. \\
 & & \left. \left. + \left(\frac{-3ikR}{k^{2}R^{2}}+\frac{3ikR}{k^{2}R^{2}}-\frac{ik^{3}R^{3}}{k^{2}R^{2}}+\frac{3ik^{3}R^{3}}{2k^{2}R^{2}}-i\frac{3k^{3}R^{3}}{6k^{2}R^{2}}+i\frac{k^{5}R^{5}}{6k^{2}R^{2}}\right)\hat{\mathbf{R}}\otimes \hat{\mathbf{R}}\right]\right\}\\\end{eqnarray*}
\begin{eqnarray*}
 & = & \lim_{R \rightarrow 0 }\left\{\frac{1}{4\pi R}\left[\left(\frac{i}{kR}+ikR-\frac{i}{kR}-\frac{ikR}{2}-i\frac{1}{6}k^{3}R^{3}+i\frac{kR}{6}\right)\mathbb{I}+ \right. \right. \\
 & & \left. \left.+ \left(\frac{-3i}{kR}+\frac{3i}{kR}-ikR+\frac{3ikR}{2}-i\frac{3kR}{6}+i\frac{k^{3}R^{3}}{6}\right)\hat{\mathbf{R}}\otimes \hat{\mathbf{R}}\right]\right\}
\end{eqnarray*}
\begin{eqnarray*}
 &  & \lim_{R \rightarrow 0 }\left\{\frac{1}{4\pi}\left[\left(\frac{i}{kR^2}-\frac{i}{kR^2}+ik-\frac{ik}{2}-i\frac{1}{6}k^{3}R^{2}+i\frac{k}{6}\right)\mathbb{I}\right.\right. \\ 
 & & \left.\left. +\left(\frac{-3i}{kR^{2}}+\frac{3i}{kR^{2}}-ik+\frac{3ik}{2}-i\frac{3k}{6}+i\frac{k^{3}R^{2}}{6}\right)\hat{\mathbf{R}}\otimes \hat{\mathbf{R}}\right]\right\}\\
  \end{eqnarray*}
 \begin{eqnarray}
 \Im\left\{ \mathbf{G}\left(\mathbf{r}_{0},\mathbf{r}_{0}\right)\right\} & = & \frac{k}{4\pi}\frac{2}{3}\mathbb{I}.
\end{eqnarray}

\section{Evaluation of the radiated power of a dipole in free space \label{App:eval_green_source}}

According with Poynting's theorem \cite{Jackson_CE}, the radiated power of a time harmonic current distribution is identical to the rate of energy dissipation:

\begin{equation}
\frac{dW}{dt}=-\frac{1}{2}\int_V^{} \Re\left\{\mathbf{j}^{*}\cdot\mathbf{E}\right\} dV.
\end{equation}
where $\mathbf{j}$ is not the total current density but instead represents the source current,  $\mathbf{j}_{s}$, that generates the fields or the loss currents,  $\mathbf{j}_{c}$, corresponding to the thermal losses and $V$ is the volume of the current distribution. If we consider the source like the one described in Eq. (\ref{eq:EM_source}), we obtain:


\begin{equation}
\frac{dW}{dt}=\frac{\omega}{2} \Im\left\{\boldsymbol{\mu}^{*}\cdot\mathbf{E}\left(\mathbf{r}_{0}\right)\right\},\label{eq:Ap1_Radiated_power}
\end{equation}
being $\mathbf{E}\left(\mathbf{r}_{0}\right)$ the electric field at dipole position at $\mathbf{r}_0$.\\
Considering the electric field (see chapter (\ref{sec:chap2_dyadic_green_function})) at point $\mathbf{r}$ as:

\begin{equation}
\mathbf{E}\left(\mathbf{r}\right) = \omega^2 \mu_{0} \mu \mathbf{G}\left(\mathbf{r},\mathbf{r}',\omega\right)\boldsymbol{\mu},
\end{equation}
Eq. (\ref{eq:Ap1_Radiated_power}) can be written as:

\begin{equation}
\frac{dW}{dt}=\frac{\omega^{3}\left|\boldsymbol{\mu}\right|^{2}}{2c^{2}\epsilon_{0}}\left[\mathbf{n}_{\mu}\cdot\Im\left\{ \mathbf{G}\left(\mathbf{r},\mathbf{r},\omega\right)\right\} \cdot\mathbf{n}_{\mu}\right],
\end{equation}
with $\mathbf{n}_{\mu}$ being the unit vector in the direction of the dipole moment. In Appendix (\ref{App:green_at_origin}) we show  how to evaluate the Green tensor at emitter position. Due to the product between $\boldsymbol{\mu}$ and $\mathbf{E}$, is necessary to evaluate the component of $\mathbf{E}$ in the direction of $\boldsymbol{\mu}$. If we place $\boldsymbol{\mu}$ in the $z-$direction (see Fig. (\ref{fig:app_axis})), $\boldsymbol{\mu} = \left|\boldsymbol{\mu}\right|\mathbf{n}_{z}$, the $z$ component of the electric field is given by:

\begin{equation}
E_z = \frac{\left|\boldsymbol{\mu}\right|}{4\pi\epsilon_{0}}\frac{e^{ikR}}{R}\left[k^2\sin^2 \theta + \frac{1}{R^2}\left(3\cos^2 \theta -1 \right) - \frac{ik}{R}\left(3 \cos^2 \theta -1\right) \right].
\end{equation}

\begin{figure}[h]
\begin{centering}
\includegraphics[width=6cm]{appendices/axis.png} 
\par\end{centering}
\caption{Schematic representation of the dipole and the corresponding field in a point in spherical coordinates, when the dipole points along the $z-$ axis.\label{fig:app_axis}}
\end{figure}
For the evaluation of the electric field at the origin of the dipole, the exponential term is expanded:

\begin{equation}
e^{ikR} = 1 + ikR - \frac{1}{2}\left(kR\right)^2  -i\frac{1}{6}\left(kR\right)^3 + \mathcal{O}^4,
\end{equation}

thus obtaining for the case where  $R \rightarrow 0$:
\begin{eqnarray}
\frac{dW}{dt} & = & \lim_{R\rightarrow0}\frac{\omega}{2}\Im\left\{ \boldsymbol{\mu}^{*}\cdot\frac{\left|\boldsymbol{\mu}\right|}{4\pi\epsilon_{0}}\frac{1}{R}\left[1+ikR+\frac{1}{2}\left(ikR\right)^{2}+\frac{1}{6}\left(ikR\right)^{3}\right]\right. \nonumber\\
 &  & \left.\left[k^{2}\sin^{2}\theta+\frac{1}{R^{2}}\left(3\cos^{2}\theta-1\right)-\frac{ik}{R}\left(3\cos^{2}\theta-1\right)\right]\right\}.
\end{eqnarray}

Considering that $\boldsymbol{\mu}$ is a real number, and only present the imaginary elements of the equation. We finally obtain:

\begin{eqnarray}
\frac{dW}{dt} & = & \frac{\left|\boldsymbol{\mu}\right|^2}{4\pi} \frac{\omega}{\epsilon_{0}}\frac{k^3}{3} \nonumber\\
 & = & \frac{\left|\boldsymbol{\mu}\right|^2}{4\pi\epsilon_{0}}\frac{k^4c}{3}.
\end{eqnarray}

%For the limit case where we obtain:
%
%\begin{equation}
%\frac{dW}{dt}=\lim_{R\rightarrow0}\frac{\omega}{2}\left|\boldsymbol{\mu}\right|\Im\left\{ E_{z}\right\} =\frac{\omega\left|\boldsymbol{\mu}\right|^{2}}{8\pi\epsilon_{0}}\lim_{R\rightarrow0}\left\{ \frac{2}{3}k^{3}+R^{2}\left(...\right)+...\right\} =\frac{\left|\boldsymbol{\mu}\right|^{2}}{12\pi}\frac{\omega}{\epsilon_{0}}k^{3},
%\end{equation}
%
%obtaining this way the radiated power by 

\section{Dipole radiation in inhomogeneous environments\label{App1:eval_dip_inh_env}}

Let us consider an emitter $\boldsymbol{\mu}$ located at $\mathbf{r} = \mathbf{r}_{0}$ surrounded by $\mathcal{N}$ dipolar particles $\mathbf{p}_n$ located at $\mathbf{r} = \mathbf{r}_{n}$ ($n = 1,...,\mathcal{N}$).

Considering the superposition principle, the total field at position $\mathbf{r}$ is given by:

\begin{equation}
\mathbf{E}\left(\mathbf{r}\right) = \frac{k^2}{\epsilon_{0}}\left[\mathbf{G}\left(\mathbf{r},\mathbf{r}_{0}\right)\cdot \boldsymbol{\mu} + \sum_{n=1}^{\mathcal{N}}\mathbf{G}\left(\mathbf{r},\mathbf{r}_{n}\right)\cdot \mathbf{p}_{n} \right].
\end{equation}
In far field, \textit{i.e.}, making the asymptotic expansion of the Green tensor $R = \left|\mathbf{r}\right| \gg \left|\mathbf{r}_{n}\right|$,  and $\mathbf{k}$ the wave vector defined by $\mathbf{k}=k \mathbf{u}_{r}$ with $\mathbf{u}_{r}=\frac{\mathbf{r}}{\left|\mathbf{r}\right|}$, the electric and magnetic field can be written as:

\begin{eqnarray}
\mathbf{E}_{n}\left(\mathbf{r}\right) & = & \frac{k^{2}}{\epsilon_{0}}\mathbf{G}\left(\mathbf{r},\mathbf{r}_{n}\right)\cdot\mathbf{p}_{n}\approx\frac{e^{ikR}}{4\pi R\epsilon_{0}}\left\{ \left(\mathbf{k}\times\mathbf{p}_{n}\right)\times\mathbf{k}\right\} e^{-i\mathbf{k}\cdot\mathbf{r}_{n}}\\
\mathbf{H}_{n}\left(\mathbf{r}\right) & = & \frac{1}{k}\sqrt{\frac{\epsilon_{0}}{\mu_{0}}}\left(\mathbf{k}\times\mathbf{E}_{n}\right)=\frac{1}{kZ}\left(\mathbf{k}\times\mathbf{E}_{n}\right),
\end{eqnarray}
where $Z=\sqrt{\mu_{0}/\epsilon_{0}}$ is the impedance.

The total power emitted will be given by:

\begin{eqnarray}
P & = & \frac{1}{2k}\int{\left(\mathbf{E}\times\mathbf{H}^{*}\right)}\cdot\left(r^{2}\sin\left(\theta\right)\right) d\theta d\phi \\
& = & \frac{1}{2Z}\int{\left|\mathbf{E}\right|^2\left[r^2\sin\left(\theta\right)\right]d\theta d\phi} \\
& = & \frac{k^{2}}{2Z\left(4\pi\epsilon_{0}\right)^{2}}\sum_{m,n}\int\left(k^{2}\mathbf{p}_{m}\cdot\mathbf{p}_{n}^{*} - \right. \nonumber\\
& & \left. -\left(\mathbf{p}_{m}\cdot\mathbf{k}\right)\left(\mathbf{p}_{n}^{*}\cdot\mathbf{k}\right)\right) e^{i\mathbf{k}\cdot\left( \mathbf{r}_{n}-\mathbf{r}_{m}\right)}\sin\left(\theta\right)d\theta d\phi.
\end{eqnarray}
The integral can be done as follows: First we consider $\theta$ as the angle measured from $\mathbf{r}_{nm}=\mathbf{r}_{n}-\mathbf{r}_{m}$. Integrating in $\phi$:

\begin{eqnarray}
& &\int\left(k^2\mathbf{p}_{m}\cdot\mathbf{p}_{n}^{*}-\left(\mathbf{p}_{m}\cdot\mathbf{k}\right)\left(\mathbf{p}_{n}^{*}\cdot\mathbf{k}\right)\right)e^{i\mathbf{k}\cdot\left(\mathbf{r}_{n}-\mathbf{r}_{m}\right)}\sin\left(\theta\right)d\theta d\phi \\
& = & k^2 2\pi \mathbf{p}_{m}\cdot\mathbf{p}_{n}^{*}\int e^{ikr_{nm}\cos\left(\theta\right)\sin\left(\theta\right)d\theta}- \nonumber \\
& & -k^2 \pi \left[ \mathbf{p}_{m} \cdot \mathbf{p}_{n}^{*} - \frac{\left(\mathbf{p}_{m}\cdot\mathbf{r}_{nm}\right)\left(\mathbf{p}_{n}^{*}\cdot\mathbf{r}_{nm}\right)}{r^{2}_{nm}}\right] \nonumber\\
& & \int \sin^{2}\left(\theta\right)e^{ikr_{nm} \cos \left(\theta\right)}\sin \left(\theta\right) d\theta \nonumber \\
& & -k^{2} 2 \pi \frac{\left(\mathbf{p}_{m}\cdot\mathbf{r}_{nm}\right)\left(\mathbf{p}_{n}^{*}\cdot\mathbf{r}_{nm}\right)}{r^{2}_{nm}} \nonumber \\
& & - \int \cos^{2}\left(\theta\right)e^{ikr_{nm} \cos \left(\theta\right)}\sin \left(\theta\right) d\theta.
\end{eqnarray}
Thus, it can be shown that the total radiated power is given by:

\begin{equation}
P = \frac{k^3}{2}\frac{1}{\epsilon_{0}\sqrt{\epsilon_{0} \mu_{0}}} \sum_{m,n}\Re \left\{ \mathbf{p}_{m}^{*}\cdot \Im \left[ \mathbf{G} \left(\mathbf{r}_{m},\mathbf{r}_{n}\right)\right]\cdot \mathbf{p}_{n} \right\}.
\end{equation}




